\documentclass[11pt]{article}

\usepackage[T1]{fontenc}
\usepackage[utf8]{inputenc}
\usepackage{lmodern}
\usepackage[margin=0.92in]{geometry}
\usepackage{amsmath,amssymb,amsthm,mathtools}
\usepackage{microtype}
\usepackage{enumitem}
\usepackage{xcolor}
\usepackage[hidelinks]{hyperref}
\usepackage[nameinlink,noabbrev]{cleveref}

\setlist[itemize]{leftmargin=1.6em,itemsep=0.25em,topsep=0.35em}
\setlist[enumerate]{leftmargin=1.8em,itemsep=0.25em,topsep=0.35em}
\setlength{\parindent}{0pt}
\setlength{\parskip}{0.55em}
\allowdisplaybreaks

\newtheorem{theorem}{Theorem}[section]
\newtheorem{proposition}[theorem]{Proposition}
\newtheorem{lemma}[theorem]{Lemma}
\newtheorem{corollary}[theorem]{Corollary}
\theoremstyle{definition}
\newtheorem{definition}[theorem]{Definition}
\newtheorem{remark}[theorem]{Remark}

\newcommand{\cU}{\mathcal U}
\newcommand{\cW}{\mathcal W}
\newcommand{\cK}{\mathcal K}
\newcommand{\N}{\mathbb N}
\newcommand{\onto}{\twoheadrightarrow}
\newcommand{\isom}{\simeq}
\newcommand{\wstar}{w^*}
\newcommand{\SpUK}{\operatorname{Sp}(\cU^{-1}\circ\cK)}

\title{\textbf{The Rosenthal--Odell Quotient Problem}\\
\large Audit of a Partial Proof, Literature Status, and Sharpened Reductions}
\author{}
\date{22 June 2026}

\begin{document}
\maketitle

\begin{abstract}
Gasparis records two closely related questions: whether every quotient of an
$\ell_1$-predual is $c_0$-saturated, and, in Odell's special case, whether every
quotient of $C(K)$ is $c_0$-saturated when $K$ is countable compact.  The phrase
``$\ell_1$-predual'' has two conventions in the literature.  Under the strict
convention, meaning that the dual is \emph{isometric} to $\ell_1$, I found no
published resolution through 22 June 2026.  Under the broader convention,
meaning only that the dual is \emph{isomorphic} to $\ell_1$, Argyros--Motakis
provide a counterexample.

The supplied complemented-subspace theorem is correct.  It is, however,
subsumed by a stronger 2017 theorem of Galego--Gonz\'alez--Pello: every quotient
$X$ of a real strict $\ell_1$-predual satisfies
\[
   \cU(X,Z)=\cK(X,Z) \qquad\text{for every Banach space } Z,
\]
where $\cU$ and $\cK$ denote the unconditionally converging and compact
operators.  This yields several additional consequences: every
infinite-dimensional quotient of $X$ contains a complemented copy of $c_0$;
$X$ has no infinite-dimensional reflexive quotient; and $X$ is
superprojective.  The remaining saturation problem is reduced exactly to the
opposite projection property:
\[
  X \text{ is } c_0\text{-saturated}
  \quad\Longleftrightarrow\quad
  X \text{ is subprojective}.
\]
The report also proves an iterative $c_0$-splitting reduction, an extension
obstruction based on the Josefson--Nissenzweig theorem, and a general
$(WU)$-trichotomy.  These sharply constrain a possible counterexample but do
not close the strict problem.
\end{abstract}

\section{The question and the convention issue}

A Banach space is $c_0$-saturated if every closed infinite-dimensional
subspace contains an isomorphic copy of $c_0$.  Gasparis asks whether every
quotient of an $\ell_1$-predual is $c_0$-saturated and attributes to Odell the
special case of quotients of $C(K)$ for countable compact metric $K$
\cite[p.~492]{Gasparis2009}.  The same problem appears in Alspach's work,
Rosenthal's survey, and Argyros--Kanellopoulos
\cite{Alspach1993,Rosenthal2003,ArgyrosKanellopoulos2005}.

The word ``predual'' is not used uniformly:

\begin{definition}
A \emph{strict $\ell_1$-predual} is a Banach space $E$ such that $E^*$ is
isometric to $\ell_1$ (or to $\ell_1(\Gamma)$ in the nonseparable notation).
A \emph{broad, or isomorphic, $\ell_1$-predual} is a Banach space $E$ such that
$E^*$ is merely isomorphic to $\ell_1$.
\end{definition}

The strict interpretation is the one relevant to the source problem.  One
clear indication is the standard assertion used there that $\ell_1$-preduals
are $c_0$-saturated, which is an isometric $L_1$-predual phenomenon.  Alspach
also explicitly adopts the isometric convention in his paper on quotients of
ordinal $C$-spaces \cite{Alspach1993}.  By contrast, Argyros--Motakis explicitly
define an $\ell_1$-predual to mean a space whose dual is only isomorphic to
$\ell_1$ \cite[p.~36]{ArgyrosMotakis2016}.

Throughout the strict part of this report, scalars are real.  This is the
setting in which the Johnson--Zippin property-$(V)$ theorem is invoked in the
supplied proof and in the cited quantitative formulation
\cite[Theorem~4.2]{Krulisova2015}.

\section{Audit of the supplied partial proof}

The supplied theorem is the following.

\begin{theorem}[Supplied partial result]\label{thm:supplied}
Let $E$ be a real Banach space such that $E^*$ is isometric to $\ell_1$, and
let $X$ be a quotient of $E$.  Every complemented infinite-dimensional
subspace of $X$ contains an isomorphic copy of $c_0$.

Consequently, the same conclusion holds when $X$ is a quotient of $C(K)$ and
$K$ is countable compact metric.
\end{theorem}

\subsection*{Verdict}
The proof is correct.  Its logical chain is:
\begin{enumerate}
\item strict real $\ell_1$-preduals have Pe\l czy\'nski's property $(V)$;
\item property $(V)$ passes to quotients;
\item an operator is unconditionally converging exactly when it fixes no copy
      of $c_0$;
\item a projection onto a $c_0$-free range is therefore weakly compact;
\item the identity on that range is weakly compact, so the range is reflexive;
\item the dual of the range embeds into $E^*=\ell_1$, hence has the Schur
      property; a reflexive Schur space is finite-dimensional.
\end{enumerate}

There is one bibliographic correction: the relevant result in
Kruli\v{s}ov\'a's paper is Theorem~4.2, not Theorem~4.5
\cite[Theorem~4.2]{Krulisova2015}.  There is also a shorter final step, given
in \cref{cor:complemented}: the projection is not only weakly compact but
compact.

For completeness, the quotient inheritance used in the proof is sound.  If
$Q:E\onto X$ and $T:X\to Z$ is unconditionally converging, then $TQ$ is
unconditionally converging.  If $TQ$ is weakly compact, the open mapping
theorem gives a constant $C$ with $B_X\subset Q(CB_E)$, and hence
$T(B_X)\subset TQ(CB_E)$ is relatively weakly compact.

\section{The stronger known operator theorem}

The main improvement comes from combining property $(V)$ with the Schur
property of the dual.  Galego--Gonz\'alez--Pello package this as the class
$\SpUK$ and prove both that strict $\ell_1$-preduals belong to it and that the
class is stable under quotients \cite[Corollary~3.5 and
Proposition~3.6]{GalegoGonzalezPello2017}.

\begin{theorem}[Compactness of unconditionally converging operators]
\label{thm:UK}
Let $E$ be a real Banach space with $E^*$ isometric to $\ell_1$, and let
$Q:E\onto X$.  Then, for every Banach space $Z$,
\[
   \cU(X,Z)=\cK(X,Z).
\]
Equivalently, every noncompact operator $T:X\to Z$ fixes an isomorphic copy of
$c_0$.
\end{theorem}

\begin{proof}
The quotient $X$ has property $(V)$.  Moreover, $Q^*:X^*\to E^*$ is an
isomorphic embedding, so $X^*$ is a subspace of a Schur space and therefore
has the Schur property.

Let $T:X\to Z$ be unconditionally converging.  Property $(V)$ makes $T$
weakly compact.  By Gantmacher's theorem, $T^*:Z^*\to X^*$ is weakly
compact.  Relatively weakly compact subsets of a Schur space are relatively
norm compact, so $T^*$ is compact.  Schauder's theorem then implies that $T$
is compact.  The reverse inclusion $\cK(X,Z)\subset\cU(X,Z)$ is standard.
\end{proof}

This direct proof also explains why the word \emph{isometric} matters: the
Schur property survives in a subspace of $\ell_1$, but property $(V)$ can fail
for a broad $\ell_1$-predual.  Galego--Gonz\'alez--Pello explicitly note that
``dual isometric'' cannot be replaced by ``dual isomorphic'' in their
corollary \cite[Corollary~3.5]{GalegoGonzalezPello2017}.

\begin{corollary}[Complemented subspaces]\label{cor:complemented}
Under the hypotheses of \cref{thm:UK}, every complemented
infinite-dimensional subspace of $X$ contains $c_0$.  In fact, if $P:X\to X$
is a projection whose range contains no $c_0$, then $P$ has finite rank.
\end{corollary}

\begin{proof}
If $P(X)$ contains no $c_0$, then $P$ fixes no copy of $c_0$, so $P$ is
unconditionally converging.  By \cref{thm:UK}, $P$ is compact.  Since $P$ is
the identity on $P(X)$, the unit ball of $P(X)$ is norm compact, and $P(X)$
is finite-dimensional.
\end{proof}

Thus the supplied argument is not merely right; its conclusion is a special
case of a stronger operator ideal statement already in the literature.

\begin{corollary}[Infinite-dimensional quotients]\label{cor:quotients}
Let $R:X\onto W$ be a quotient map with $W$ infinite-dimensional.  Then
$R$ fixes a copy $A\isom c_0$.  The copy $R(A)\isom c_0$ is complemented in
$W$, and $A$ may be chosen complemented in $X$.  Consequently, $X$ has no
infinite-dimensional reflexive quotient.
\end{corollary}

\begin{proof}
A compact surjection has finite-dimensional range, so $R$ is noncompact.
By \cref{thm:UK}, $R$ is not unconditionally converging and therefore fixes
a copy $A\isom c_0$.  The spaces $E$, $X$, and $W$ are separable because
$E^*$ is separable.  Sobczyk's theorem complements $R(A)$ in $W$.  If
$P:W\to R(A)$ is a projection and $S:R(A)\to A$ is the inverse of $R|_A$,
then $SPR:X\to A$ is a projection onto $A$.

If $W$ were infinite-dimensional reflexive, it could not contain $c_0$,
contradicting the first part.
\end{proof}

\begin{corollary}[Superprojectivity]\label{cor:super}
Every quotient $X$ of a strict real $\ell_1$-predual is superprojective.
\end{corollary}

\begin{proof}
This is Corollary~3.4 together with Corollary~3.5 and Proposition~3.6 of
\cite{GalegoGonzalezPello2017}.  It also follows directly from
\cref{cor:quotients}: for an infinite-codimensional subspace $M\subset X$,
the quotient map $X\to X/M$ fixes a copy of $c_0$ whose image is complemented
in the separable quotient.  Pulling back a complementary summand gives an
infinite-codimensional complemented subspace of $X$ containing $M$.
\end{proof}

\section{An exact reduction: saturation versus subprojectivity}

Recall that $X$ is \emph{subprojective} if every infinite-dimensional
subspace of $X$ contains an infinite-dimensional subspace complemented in
$X$.

\begin{theorem}[Exact reduction]\label{thm:equiv}
Let $X$ be a quotient of a real strict $\ell_1$-predual.  Then
\[
   X \text{ is }c_0\text{-saturated}
   \quad\Longleftrightarrow\quad
   X \text{ is subprojective}.
\]
The same equivalence holds for every quotient of $C(K)$ with $K$ countable
compact metric.
\end{theorem}

\begin{proof}
Suppose first that $X$ is $c_0$-saturated.  If $Y\subset X$ is
infinite-dimensional, choose $C\subset Y$ with $C\isom c_0$.  Since $X$ is
separable, Sobczyk's theorem makes $C$ complemented in $X$.  Thus $X$ is
subprojective.  This implication is also a special case of a theorem of
D\'iaz--Fern\'andez, as recorded in
\cite[Proposition~2.7]{GalegoGonzalezPello2017}.

Conversely, suppose $X$ is subprojective and let $Y\subset X$ be
infinite-dimensional.  There is an infinite-dimensional $Z\subset Y$
complemented in $X$.  By \cref{cor:complemented}, $Z$, and hence $Y$,
contains $c_0$.
\end{proof}

Accordingly, the strict Rosenthal--Odell problem can be restated as follows.

\begin{quote}
\textbf{Equivalent strict problem.} Is every quotient of a strict
$\ell_1$-predual subprojective?  In particular, is every quotient of
$C(K)$ subprojective when $K$ is countable compact?
\end{quote}

This reformulation is useful because the \emph{known} direction is the dual
one: such quotients are superprojective by \cref{cor:super}.  The 2025 work of
Gonz\'alez--Pello emphasizes that subprojectivity passes to subspaces,
superprojectivity passes to quotients, and neither property is a general
three-space property \cite{GonzalezPello2025}.  Thus the missing implication
is not a formal permanence fact.

\section{Literature search: what is solved and what remains}

\subsection{The strict/isometric version}

I found no published proof or counterexample for the strict version through
22 June 2026.  The following checkpoints support that assessment.

\begin{itemize}
\item Alspach wrote in 1992 that it was unknown whether quotients of
      $C(\alpha)$ must be $c_0$-saturated \cite{Alspach1993}.
\item Odell proved that quotients of Schreier's space are $c_0$-saturated,
      and listed the higher Schreier and countable-$C(K)$ questions as open
      \cite{Odell1992}.
\item Rosenthal's 2003 survey states that the countable-$C(K)$ quotient
      problem is unknown and notes that even the possible appearance of
      $\ell_2$ inside a quotient of $C(\omega^\omega)$ was unknown
      \cite{Rosenthal2003}.
\item Argyros--Kanellopoulos again call the problem well known and open
      \cite[Remark~4]{ArgyrosKanellopoulos2005}.
\item Gasparis repeats the question in the source paper
      \cite[p.~492]{Gasparis2009}.
\item A 2020 MathOverflow formulation by Argyros still describes it as a
      well-known problem \cite{MathOverflow2020}.
\item The 2017 operator theorem gives the strong conclusions in
      \cref{thm:UK,cor:super}, but it proves superprojectivity rather than the
      subprojectivity required by \cref{thm:equiv}.
\item The 2025 three-space paper develops new criteria for subprojectivity
      and superprojectivity but does not provide the missing quotient
      permanence \cite{GonzalezPello2025}.
\end{itemize}

This is a report of the search, not a claim that no inaccessible manuscript
exists.  The precise conclusion is that no published resolution was located
in the supplied papers, arXiv searches, journal searches, bibliographic
forward searches, or the recent subprojectivity literature.

\subsection{The broad/isomorphic version has a counterexample}

Under the convention $E^*\isom\ell_1$, the answer is negative.
Argyros--Motakis construct a hereditarily indecomposable
$\mathcal L_\infty$-space $\mathfrak X_{\mathrm{nr}}$ containing no $c_0$, no
$\ell_1$, and no reflexive subspace \cite{ArgyrosMotakis2016}.  Their
Theorem~8.3 yields infinite-dimensional spaces
\[
    X_1\onto X_2\onto X_3
\]
where all three are broad $\ell_1$-preduals with the scalar-plus-compact
property, $X_1$ and $X_3$ are reflexively saturated, and the middle space is
the no-reflexive-subspace construction.  The preceding discussion identifies
$X_2$ with $\mathfrak X_{\mathrm{nr}}$ (or the corresponding version), and
Corollary~3.12 proves that it contains no $c_0$.

Hence $X_1\onto X_2$ is a quotient of a broad $\ell_1$-predual that is not
$c_0$-saturated.  This fully settles the broad reading, but not the source
question.  Indeed, it cannot be converted naively into a strict example: if
$X_1$ were a quotient of a strict $\ell_1$-predual, then the quotient map onto
the $c_0$-free infinite-dimensional space $X_2$ would be unconditionally
converging and therefore compact by \cref{thm:UK}, an impossibility.

\section{Further structural reductions}

The remainder records additional deductions that were not present in the
supplied packet.  They do not solve the problem, but they describe the shape
of any strict counterexample.

\subsection{A residual $c_0$-splitting lemma}

\begin{lemma}[Transverse $c_0$ splitting]\label{lem:transverse}
Let $X$ be a quotient of a real strict $\ell_1$-predual, let
$Y\subset X$ contain no copy of $c_0$, and let $q:X\to X/Y$ be the quotient
map.  If $A\subset X$ is isomorphic to $c_0$, then:
\begin{enumerate}
\item $q|_A$ is not compact;
\item there is a further $B\subset A$, $B\isom c_0$, such that $q|_B$ is an
      isomorphism;
\item $q(B)$ is complemented in $X/Y$, and $B$ is complemented in $X$.
\end{enumerate}
\end{lemma}

\begin{proof}
Choose an isomorphism $U:c_0\to A$ and put $a_n=Ue_n$.  If $q|_A$ were
compact, then $\|q(a_n)\|\to0$, because $(a_n)$ is weakly null.  Passing to a
subsequence, choose $y_n\in Y$ so that
$\sum_n\|a_n-y_n\|$ is small enough for the principle of small perturbations.
Then $(y_n)$ is equivalent to a subsequence of the $c_0$ basis, contradicting
the assumption on $Y$.  This proves (1).

Since $A\isom c_0$, every unconditionally converging operator from $A$ is
compact (apply \cref{thm:UK} to $c_0$, or use property $(V)$ and
$c_0^*=\ell_1$).  Thus the noncompact operator $q|_A$ is not unconditionally
converging and fixes a copy $B\isom c_0$, proving (2).

The quotient $X/Y$ is separable, so Sobczyk's theorem complements $q(B)$.
If $P:X/Y\to q(B)$ is a projection and $S=(q|_B)^{-1}$, then $SPq$ is a
projection from $X$ onto $B$.
\end{proof}

\begin{corollary}[Iterated peeling]\label{cor:peeling}
If a strict counterexample contains an infinite-dimensional $c_0$-free
subspace $Y\subset X$, then for every $n$ there are complemented copies
$B_1,\ldots,B_n\isom c_0$ and a complemented subspace $X_n\subset X$ such
that
\[
    X = B_1\oplus\cdots\oplus B_n\oplus X_n,
    \qquad Y\subset X_n,
\]
and $X_n$ is again a quotient of a strict $\ell_1$-predual.  Moreover,
$X_n/Y$ remains infinite-dimensional.
\end{corollary}

\begin{proof}
The identity on the infinite-dimensional $X$ is noncompact, so
\cref{thm:UK} gives a copy of $c_0$ in $X$.  Apply
\cref{lem:transverse}.  The resulting projection onto $B_1$ has a kernel
$X_1$ containing $Y$, and $X_1$ is a quotient of $X$ via the complementary
projection; hence it remains in the class of \cref{thm:UK}.  If $X_1/Y$ were
finite-dimensional, then $Y$ would be complemented in $X_1$, contradicting
\cref{cor:complemented}.  Iterate.
\end{proof}

This produces arbitrarily long $c_0$ decompositions transverse to a bad
subspace.  It does not by itself produce a $c_0$ sequence in $Y$, because the
norms of the successive projections need not be uniformly controlled.  That
loss of quantitative control is a concrete obstruction to turning the
peeling argument into a proof.

\subsection{A nonextendability phenomenon}

\begin{proposition}[Unconditionally converging extensions must fail]
\label{prop:extension}
Let $Y\subset X$ be infinite-dimensional and $c_0$-free, where $X$ is as in
\cref{thm:UK}.  There exists a bounded, noncompact, unconditionally
converging operator $T:Y\to c_0$.  Every bounded extension
$\widetilde T:X\to c_0$ of $T$ fixes a copy of $c_0$; in particular, no such
extension is unconditionally converging.
\end{proposition}

\begin{proof}
By the Josefson--Nissenzweig theorem, choose a normalized weak-star null
sequence $(y_n^*)\subset Y^*$.  Define
\[
       T(y)=(y_n^*(y))_{n=1}^\infty\in c_0.
\]
Then $T$ is bounded.  It is not compact: $T^*e_n^*=y_n^*$ has norm one,
whereas compactness would force this weak-star null sequence to have a
norm-null subsequence.  Since $Y$ contains no $c_0$, every operator defined
on $Y$ is unconditionally converging.

The separability of $X$ and Sobczyk's extension theorem provide a bounded
extension $\widetilde T:X\to c_0$.  If $\widetilde T$ were unconditionally
converging, \cref{thm:UK} would make it compact, and then its restriction
$T$ would be compact, a contradiction.  Hence every extension fixes $c_0$.
\end{proof}

This proposition isolates the hereditary gap.  On the whole space,
unconditionally converging means compact; on a hypothetical $c_0$-free
subspace, there are automatically noncompact unconditionally converging
operators.  Every extension is forced to acquire a $c_0$ copy outside the
subspace.

\subsection{The $(WU)$ trichotomy}

Odell calls a Banach space $X$ a $(WU)$ space if every normalized weakly null
sequence in $X$ has an unconditional subsequence.  Every quotient of a space
with a shrinking unconditional finite-dimensional decomposition has $(WU)$
\cite[Theorem~A]{Odell1992}.

\begin{theorem}[General $(WU)$ trichotomy]\label{thm:WU}
Let $X$ have property $(WU)$.  Every infinite-dimensional subspace
$Y\subset X$ contains a copy of $\ell_1$, a copy of $c_0$, or an
infinite-dimensional reflexive subspace.
\end{theorem}

\begin{proof}
Assume that $Y$ contains neither $\ell_1$ nor $c_0$.  If $Y$ is reflexive,
we are done.  Otherwise, by Eberlein--\v{S}mulian choose a bounded sequence
with no weakly convergent subsequence.  Rosenthal's $\ell_1$ theorem gives a
weakly Cauchy subsequence $(y_n)$, which is not weakly convergent and hence
not norm Cauchy.  Select
$n_1<m_1<n_2<m_2<\cdots$ and $\delta>0$ with
$\|y_{m_j}-y_{n_j}\|\ge\delta$, and set
\[
   z_j=\frac{y_{m_j}-y_{n_j}}{\|y_{m_j}-y_{n_j}\|}.
\]
Then $(z_j)$ is normalized and weakly null.  Property $(WU)$ gives an
unconditional basic subsequence.  Its closed span contains neither $\ell_1$
nor $c_0$.  By James's criteria, the unconditional basis is both shrinking
and boundedly complete, so its span is reflexive
\cite{James1950,Rosenthal1974}.
\end{proof}

\begin{corollary}\label{cor:WU}
If $X$ has $(WU)$ and contains no $\ell_1$, then
\[
 X \text{ is }c_0\text{-saturated}
 \quad\Longleftrightarrow\quad
 X \text{ contains no infinite-dimensional reflexive subspace}.
\]
In particular, this applies to quotients of spaces with shrinking
unconditional FDDs.
\end{corollary}

For a quotient $X$ of countable $C(K)$, every subspace $Y\subset X$ has
separable dual: $X^*$ embeds into $C(K)^*=\ell_1(K)$ and the restriction
map $X^*\onto Y^*$ is onto.  Hence no such $Y$ contains $\ell_1$.  For the
first hard level $C(\omega^\omega)$, Odell also records Schlumprecht's
observation that a normalized weakly null sequence with an $\ell_1$
spreading model generates a subspace containing $c_0$
\cite[Section~3]{Odell1992}.  Therefore a $c_0$-free, no-reflexive-subspace
counterexample at that level would have to fail $(WU)$ through a
Maurey--Rosenthal type sequence, while simultaneously avoiding every
$\ell_1$ spreading model.

\section{Why the natural closing arguments fail}

The preceding results leave several tempting but invalid shortcuts.

\begin{enumerate}
\item \textbf{The Schur property does not pass to the dual of a subspace.}
      Although $X^*\subset\ell_1$ is Schur, for $Y\subset X$ the space
      $Y^*$ is a quotient of $X^*$, and quotients of $\ell_1$ can be
      arbitrary separable Banach spaces.  Thus $Y^*$ may be reflexive or
      James-like.

\item \textbf{Superprojectivity points in the wrong direction.}
      It controls infinite-codimensional superspaces and quotients.  The
      desired conclusion is subprojectivity, which controls subspaces.
      Neither is a formal consequence of the other.

\item \textbf{Cantor--Bendixson induction has a nonclosed-range gap.}
      For countable $K$, the natural kernel
      $C_0(K\setminus K')\isom c_0$ suggests transfinite induction.  But the
      image of this kernel under an arbitrary quotient map need not be
      closed.  The closure of a nonclosed operator range from $c_0$ can hide
      a $c_0$-free separable space, so the exact sequence alone does not
      preserve saturation.

\item \textbf{The broad counterexample cannot simply be renormed.}
      The strict hypothesis supplies property $(V)$ and hence
      $\cU=\cK$ after passing to quotients.  The Argyros--Motakis quotient
      onto a $c_0$-free space violates exactly this operator conclusion.

\item \textbf{Iterated splitting lacks uniform constants.}
      \Cref{cor:peeling} gives arbitrarily many complemented $c_0$ pieces
      away from a bad $Y$, but the projection norms and isomorphism constants
      may diverge.  Without uniform control, no bounded $c_0$-sum projection
      or basis decomposition follows.
\end{enumerate}

These are not merely technical annoyances: they identify the point at which
each natural proof strategy requires genuinely new information.

\section{Final status}

\begin{enumerate}
\item The supplied complemented-subspace result is correct.  Its only needed
      repair is the citation to Kruli\v{s}ov\'a's Theorem~4.2.
\item A stronger known theorem gives $\cU(X,Z)=\cK(X,Z)$ for every quotient
      $X$ of a strict real $\ell_1$-predual and every $Z$.
\item Under the broad convention $X^*\isom\ell_1$, the literature contains a
      full counterexample, via Argyros--Motakis.
\item Under the strict convention intended by the source, and in particular
      for quotients of countable $C(K)$, no full proof or counterexample was
      located through 22 June 2026.
\item The strict problem is exactly the subprojectivity of these quotients.
      Any counterexample must support the transverse splitting and extension
      phenomena proved above, while avoiding the standard $(WU)$ and
      spreading-model routes to $c_0$.
\end{enumerate}

\begin{thebibliography}{99}

\bibitem{Alspach1993}
D. E. Alspach,
\emph{A $\ell_1$-predual which is not isometric to a quotient of
$C(\alpha)$},
Banach Spaces (M\'erida, 1992), Contemp. Math. \textbf{144},
American Mathematical Society, 1993, 9--14.

\bibitem{ArgyrosKanellopoulos2005}
S. A. Argyros and V. Kanellopoulos,
\emph{Determining $c_0$ in $C(K)$ spaces},
Fund. Math. \textbf{187} (2005), 61--93.

\bibitem{ArgyrosMotakis2016}
S. A. Argyros and P. Motakis,
\emph{The scalar-plus-compact property in spaces without reflexive
subspaces},
arXiv:1608.01962, 2016.

\bibitem{Diestel1984}
J. Diestel,
\emph{Sequences and Series in Banach Spaces},
Graduate Texts in Mathematics 92, Springer, 1984.

\bibitem{GalegoGonzalezPello2017}
E. M. Galego, M. Gonz\'alez, and J. Pello,
\emph{On subprojectivity and superprojectivity of Banach spaces},
Results Math. \textbf{71} (2017), 1191--1205.

\bibitem{Gasparis2009}
I. Gasparis,
\emph{New examples of $c_0$-saturated Banach spaces},
Math. Ann. \textbf{344} (2009), 491--500.

\bibitem{GonzalezPello2025}
M. Gonz\'alez and J. Pello,
\emph{On the three-space property for subprojective and superprojective
Banach spaces},
Mediterr. J. Math. \textbf{22} (2025), article 54.

\bibitem{James1950}
R. C. James,
\emph{Bases and reflexivity of Banach spaces},
Ann. of Math. (2) \textbf{52} (1950), 518--527.

\bibitem{Krulisova2015}
H. Kruli\v{s}ov\'a,
\emph{Quantification of Pe\l czy\'nski's property $(V)$},
arXiv:1509.06610, 2015.

\bibitem{MathOverflow2020}
S. A. Argyros,
\emph{Quotients of subspaces of $C(\alpha)$},
MathOverflow question 360068, 11 May 2020,
\url{https://mathoverflow.net/questions/360068/quotients-of-subspaces-of-c-alpha}.

\bibitem{Odell1992}
E. Odell,
\emph{On quotients of Banach spaces having shrinking unconditional bases},
Illinois J. Math. \textbf{36} (1992), 681--695.

\bibitem{Pelczynski1962}
A. Pe\l czy\'nski,
\emph{Banach spaces in which every weakly unconditionally converging
operator is weakly compact},
Bull. Acad. Polon. Sci. S\'er. Sci. Math. Astronom. Phys. \textbf{10}
(1962), 641--648.

\bibitem{Rosenthal1974}
H. P. Rosenthal,
\emph{A characterization of Banach spaces containing $\ell_1$},
Proc. Natl. Acad. Sci. USA \textbf{71} (1974), 2411--2413.

\bibitem{Rosenthal2003}
H. P. Rosenthal,
\emph{The Banach spaces $C(K)$},
in Handbook of the Geometry of Banach Spaces, Vol. 2,
North-Holland, 2003, 1547--1602.

\bibitem{Sobczyk1941}
A. Sobczyk,
\emph{Projection of the space $(m)$ on its subspace $(c_0)$},
Bull. Amer. Math. Soc. \textbf{47} (1941), 938--947.

\end{thebibliography}

\end{document}
